% Spatial extent: where the elementary reconstruction stops.
% Lane: O (operator bridge), modules YangMills/OS/**.
% REGIME: tricotomy; NO scores, NO desk language, NO commit chronology beyond
% the reproducibility freeze.  Freeze: Lean tree of source checkpoint
% d8187d7f9caa9ce93a58a70a382e98a239942afe, full core build 8424 jobs,
% 2481 oracle commands (see Section 6).
\documentclass[11pt]{article}
\usepackage[margin=1.1in]{geometry}
\usepackage{amsmath,amssymb,amsthm}
\usepackage{booktabs,longtable,array}
\usepackage[colorlinks=true,linkcolor=blue,urlcolor=blue,citecolor=blue]{hyperref}

\newtheorem{theorem}{Theorem}[section]
\newtheorem{lemma}[theorem]{Lemma}
\newtheorem{proposition}[theorem]{Proposition}
\newtheorem{corollary}[theorem]{Corollary}
\newtheorem{remark}[theorem]{Remark}
\theoremstyle{definition}
\newtheorem{definition}[theorem]{Definition}

\emergencystretch=3em
\tolerance=2000

\newcommand{\lean}[1]{\texttt{#1}}
\newcommand{\anchor}{d8187d7f9caa9ce93a58a70a382e98a239942afe}
\newcommand{\osline}[3]{%
  \href{https://github.com/lluiseriksson/THE-ERIKSSON-PROGRAMME/blob/\anchor/YangMills/OS/#1.lean\#L#2}{\nolinkurl{#3}}}
\newcommand{\repofile}[2]{%
  \href{https://github.com/lluiseriksson/THE-ERIKSSON-PROGRAMME/blob/\anchor/#1}{\nolinkurl{#2}}}
\newcommand{\R}{\mathbb{R}}
\newcommand{\C}{\mathbb{C}}
\newcommand{\Z}{\mathbb{Z}}
\newcommand{\SU}{\mathrm{SU}}
\newcommand{\Om}{\Omega}

\title{Where the Elementary Reconstruction Stops:\\
Spatial Coupling Breaks the Uniform Vacuum, Machine-Checked\\
\large A negative result, and the positive one that isolates it}
\author{Lluis Eriksson}
\date{}

\begin{document}
\maketitle

\begin{abstract}
Two companion developments verified an Osterwalder--Seiler reconstruction end to
end for a lattice gauge chain whose spatial slice is a single point.  Every step
of that chain begins by knowing the vacuum, and in the one-dimensional case the
vacuum is free: the normalised transfer kernel has constant row sums, so the
uniform vector is fixed, and $T\Om=\Om$ follows from normalisation alone.  This
paper asks what survives when the slice acquires spatial extent, and answers in
Lean~4 with \texttt{mathlib}.

The algebraic half survives untouched.  With time bonds only, the row sums of
the transfer kernel are constant for every spatial extent $L$, so the uniform
vacuum persists on a space of dimension $2^L$; and the single-site sign
observable is an eigenvector whose normalised eigenvalue is exactly $\tanh\beta$,
with $L$ free in the statement.

The vacuum does not survive.  Switching on a coupling between sites inside a
slice makes the spatial weight depend only on the source configuration, so it
factors out of the sum over the target and the row sums become
configuration-dependent.  We exhibit two explicit configurations of a two-site
slice with different row sums, and conclude that no constant row sum exists: the
uniform vector is not fixed, so $T\Om=\Om$ is \emph{false} for it.  The vacuum
becomes a Perron vector that row-sum normalisation no longer supplies in closed
form, and every later step of the reconstruction loses its starting point.

We state plainly what the positive half is and is not.  The decoupled system is
$L$ non-interacting copies of a two-state system, and the largest non-vacuum
eigenvalue of that tensor product is independent of $L$ for trivial reasons, so
the rate it yields is physically empty
and is recorded only because it isolates which half of the construction
survives.  \textbf{No gap for the coupled system is proved here, and none is
claimed.}  Nothing in this paper is a claim about $\SU(N)$, the continuum limit,
or the Yang--Mills mass gap.
\end{abstract}

\section{Introduction}

\subsection{A prediction, and a test of it}

The companion papers \cite{ErikssonOSChain, ErikssonQuotient} verified, for the
$\Z_2$ chain, a complete passage from a Gibbs weight to exponential decay of a
correlation function: reflection positivity, a Gelfand--Naimark--Segal space, a
transfer operator, an identification theorem, a spectral bound, and an exact
two-point function.  Both papers state in print that the next step is spatial
extent, and both predict that it is where the construction stops being easy.

This paper tests that prediction rather than repeating it.  The result is that
the prediction was right, and that it is possible to say exactly \emph{which}
step fails and why.

\subsection{Why the vacuum is the load-bearing step}

Every theorem of the one-dimensional chain begins by knowing the vacuum.  The
identification is a statement about $\langle A\Om, T^n B\Om\rangle$; the
spectral bound is a bound on $T-|\Om\rangle\langle\Om|$; the endpoint is an
expectation evaluated against $\Om$.  And in one dimension the vacuum is free:
the normalised kernel has constant row sums, so the uniform vector is fixed, and
$T\Om=\Om$ is a consequence of normalisation --- the companion paper makes a
point of the fact that it is a consequence and not a hypothesis.

If that step fails, nothing downstream has an input.  That is what makes it the
right thing to test first, and it is what this paper tests.

\subsection{What is proved}

Slices carry $L$ spins, so configurations are maps $\Z_2^{L}$ and the transfer
operator acts on a space of dimension $2^{L}$.  Write $k_\beta$ for the
Boltzmann bond weight and $Z_\beta=e^{\beta}+e^{-\beta}$.

\medskip
\noindent\textbf{Decoupled} --- kernel $K(\sigma,\tau)=\prod_j k_\beta(\sigma_j,\tau_j)$:
\begin{itemize}
\item $\sum_\tau K(\sigma,\tau) = Z_\beta^{L}$, independent of $\sigma$, for
      arbitrary $L$;
\item the single-site sign observable is an eigenvector, with normalised
      eigenvalue exactly $\tanh\beta$ --- the same number as in one dimension,
      with $L$ free.
\end{itemize}

\medskip
\noindent\textbf{Coupled} --- the same, times a weight $w_\gamma(\sigma)$
linking sites inside the slice:
\begin{itemize}
\item $\sum_\tau K_\gamma(\sigma,\tau) = w_\gamma(\sigma)\,Z_\beta^{L}$;
\item two explicit configurations of a two-site slice have different row sums
      at every $\gamma>0$;
\item hence no constant row sum exists, and $T\Om=\Om$ fails.
\end{itemize}

\subsection{What is deliberately not claimed}

This subsection is not boilerplate; the first item is the one that matters, and
it was fixed in the governing charter before any of this was written.

\begin{enumerate}
\item \textbf{The volume-independent rate is physically empty.}  The decoupled
system is $L$ copies of a two-state system with no interaction between them, and
the largest non-vacuum eigenvalue of that tensor product does not depend on $L$.  We do not
call this a volume-uniform gap, we do not headline it, and it is
\emph{not} evidence about the coupled case.  It is recorded because it isolates
which half of the construction survives.
\item \textbf{No gap for the coupled system.}  None is proved and none is
claimed.  That is the open problem.
\item \textbf{Still $\Z_2$}, one variable per site, fixed finite size, and the
spatial coupling is exhibited on a two-site slice.
\item \textbf{Nothing about $\SU(N)$, the continuum limit, or the Clay problem}
\cite{Clay}.  No reduction is claimed and none is known to the author.
\end{enumerate}

\section{The factorisation that carries both halves}

\begin{lemma}\label{lem:factor}
For $f:\{1,\dots,L\}\times\Z_2\to\R$,
$\displaystyle\sum_{\tau}\prod_j f_j(\tau_j) \;=\; \prod_j \sum_{t} f_j(t)$.
\end{lemma}

\osline{SpatialExtent}{77}{sum\_prod\_config}.  Both the positive and the
negative half below are corollaries of this one identity applied to different
$f$.

\section{The decoupled system: the algebraic half ports}

\begin{theorem}[Constant row sums at every spatial extent]\label{thm:rows}
$\displaystyle\sum_{\tau} K(\sigma,\tau) = Z_\beta^{L}$ for every $\sigma$ and
every $L$.
\end{theorem}

\begin{proof}
Lemma~\ref{lem:factor} with $f_j(t)=k_\beta(\sigma_j,t)$; each factor is
$\sum_t k_\beta(s,t)=Z_\beta$ whatever $s$ is.
\end{proof}

\osline{SpatialExtent}{99}{sum\_spatialKernel}, and as an existence statement
\osline{SpatialExtent}{232}{decoupled\_constant\_rowSum}.  This is what carries
the uniform vacuum to arbitrary spatial extent.

\begin{theorem}[A mode whose rate does not depend on the volume]\label{thm:mode}
Let $s_i(\sigma)$ be the sign of site $i$.  Then
$\displaystyle\sum_\tau K(\sigma,\tau)\, s_i(\tau)
 = \bigl(e^{\beta}-e^{-\beta}\bigr) Z_\beta^{L-1}\, s_i(\sigma)$,
so after dividing by the row sum the eigenvalue is
$\bigl(e^{\beta}-e^{-\beta}\bigr)/Z_\beta=\tanh\beta$.
\end{theorem}

\osline{SpatialExtent}{129}{spatialKernel\_siteSign} and
\osline{SpatialExtent}{160}{spatial\_rate\_eq\_tanh}.  In the formal statement
$L$ is a free variable, so the independence is visible in the statement rather
than inferred from a family of instances.

\begin{remark}[What Theorem~\ref{thm:mode} is worth]\label{rem:empty}
Almost nothing on its own, and it is important to say so.  The decoupled system
is a tensor power; its spectrum is products of one-site eigenvalues, and the
largest non-vacuum one is volume-independent for the same reason that
$1^{L-1}\lambda=\lambda$.  A reader who takes this as progress towards a
volume-uniform statement about an interacting theory has been misled, and we
would rather forestall that here than be quoted out of context later.  The value
of the theorem is comparative: it fixes what the coupled case has to be measured
against.
\end{remark}

\section{The coupled system: the vacuum breaks}

\begin{definition}
$K_\gamma(\sigma,\tau) = w_\gamma(\sigma)\,K(\sigma,\tau)$, where
$w_\gamma$ is a bond weight between two sites of the slice.
\end{definition}

The essential feature is that $w_\gamma$ depends on the \emph{source} only:
\osline{SpatialExtent}{171}{spatialWeight},
\osline{SpatialExtent}{176}{coupledKernel}.

\begin{theorem}[The row sum becomes configuration-dependent]\label{thm:coupled}
$\displaystyle\sum_\tau K_\gamma(\sigma,\tau) = w_\gamma(\sigma)\,Z_\beta^{L}$.
\end{theorem}

\osline{SpatialExtent}{181}{sum\_coupledKernel}.  The spatial weight is constant
in $\tau$, so it leaves the sum untouched and survives into the answer.

\begin{theorem}[The obstruction]\label{thm:obstruction}
For every $\gamma>0$ there are two configurations of a two-site slice with
different row sums.  Consequently there is no constant $c$ with
$\sum_\tau K_\gamma(\sigma,\tau)=c$ for all $\sigma$, and the uniform vector is
not fixed: $T\Om=\Om$ is false.
\end{theorem}

\begin{proof}
Take $\sigma$ constant and $\sigma'$ alternating.  By
Theorem~\ref{thm:coupled} the row sums are $e^{\gamma}Z_\beta^{2}$ and
$e^{-\gamma}Z_\beta^{2}$, and $e^{-\gamma}<e^{\gamma}$ for $\gamma>0$.
\end{proof}

Artifacts: \osline{SpatialExtent}{188}{cfgFlat},
\osline{SpatialExtent}{190}{cfgAlt},
\osline{SpatialExtent}{208}{coupled\_rowSums\_not\_constant},
\osline{SpatialExtent}{224}{coupled\_no\_constant\_rowSum}.

\begin{remark}[Why this ends the elementary route]
In one dimension $T\Om=\Om$ was free.  It is not a small convenience: the
identification, the spectral bound and the endpoint of the companion papers all
take the vacuum as an input, and all three are stated in terms of it.  Once the
row sums vary with the configuration, the top eigenvector is a Perron vector
that row-sum normalisation no longer supplies, and none of those statements can
be written down without first solving for it.  The obstruction is therefore not that a proof got
harder; it is that the first object in the chain stopped being available.
\end{remark}

\begin{remark}[The symmetric kernel, and what is and is not formalised]\label{rem:sym}
The coupled kernel above carries the spatial weight on the source only, which
makes the row sum transparent but is not symmetric.  The form usual in the
physics literature is the symmetrised
\[
  \widetilde K_\gamma(\sigma,\tau) \;=\; w_\gamma(\sigma)^{1/2}\, K(\sigma,\tau)\, w_\gamma(\tau)^{1/2},
\]
which is $D^{-1}K_\gamma D$ for $D=\mathrm{diag}(w_\gamma^{1/2})$ and therefore has
the same spectrum.  Three things should be kept apart.
\begin{itemize}
\item \textbf{What is proved:} for the oriented kernel, the row sums are not
constant and the uniform vector is not fixed.  That is the statement the
elementary route actually consumes, because that route obtains the vacuum
\emph{from} constant row sums.
\item \textbf{What is asserted and not formalised:} that the symmetrised kernel
does not fix the uniform vector either.  The diagonal similarity above is
standard, but neither it nor that consequence is machine-checked in this
development, and we mark it as prose rather than let it pass for a theorem.
\item \textbf{What is \emph{not} claimed at all:} that the transfer operator ceases
to exist, or that no vacuum exists.  A positive kernel on a finite set has a
Perron vector.  The obstruction is that the elementary normalisation stops
handing it to us, not that the object is absent.
\end{itemize}
A reader who suspects the obstruction is an artifact of a kernel convention is
asking the right question.  The convention decides which matrix one writes
down; it does not decide whether row-sum normalisation still produces the
vacuum, and it is the latter that the chain of \cite{ErikssonOSChain} consumed.
\end{remark}
\begin{remark}[What would have to happen next]
The open problem is a spectral bound for a transfer operator whose vacuum is not
known in closed form.  That is the genuine analytic content of the subject, it is
where the constructive literature is difficult, and nothing in the present
development touches it.  Two honest continuations exist: a Perron--Frobenius
route that produces the vacuum at spatial extent without a closed form, or a
statement that the elementary methods of this lane end here.  We do not know
which, and we decline to estimate.
\end{remark}

\section{Formalisation notes}

\paragraph{One lemma, two opposite conclusions.}  Lemma~\ref{lem:factor} proves
both that the decoupled row sums are constant and that the coupled ones are not.
The difference is entirely in which function is fed to it, which is a compact
way of exhibiting that the obstruction is structural rather than incidental.

\paragraph{A shared case split.}  Three separate proofs originally repeated the
same two-case analysis of a single site, each slightly differently, and three
of the module's failures had that single cause.  Extracting the case split and
the one-site row sum as shared lemmas removed all three at once:
\osline{SpatialExtent}{60}{fin\_two\_cases},
\osline{SpatialExtent}{65}{sum\_bond\_row}.

\paragraph{A linter warning obeyed without checking.}  A simplification step
flagged as unused was in fact load-bearing; removing it broke a proof whose
error message pointed elsewhere.  The reduction is now performed explicitly so
that it does not depend on a linter's opinion.  We record this because it is the
second occurrence in this development.

\section{Reproducibility}\label{sec:repro}

All artifacts are at source checkpoint \texttt{\anchor} on branch \texttt{main}
of \url{https://github.com/lluiseriksson/THE-ERIKSSON-PROGRAMME}.  Every link
was resolved against that commit, and checked to point at an actual declaration
line, before compilation.

\begin{center}
\begin{tabular}{ll}
\toprule
Quantity & Value \\
\midrule
Full core build & 8424 jobs, success \\
Oracle commands & 2481 \\
\quad with axiom dependencies & 2459 \\
\quad axiom-free & 22 \\
Distinct axioms across all outputs & \texttt{propext}, \texttt{Quot.sound}, \texttt{Classical.choice} \\
Occurrences of \texttt{sorryAx} & 0 \\
Project axioms & 0 \\
Errors & 0 \\
\bottomrule
\end{tabular}
\end{center}

The criteria this work had to meet --- including the requirement that the
obstruction be a proved negative with explicit witnesses, and the requirement
that the decoupled result not be presented as evidence about the coupled case
--- were fixed in
\repofile{docs/O-BRIDGE-CHARTER.md}{docs/O-BRIDGE-CHARTER.md}, Amendment~10, in
a commit preceding the existence of the module they govern.  Their verdicts are
in \repofile{docs/VERIFICATION-LEDGER.md}{docs/VERIFICATION-LEDGER.md},
Addendum~520.  The oracle script is
\repofile{oracle\_check.lean}{oracle\_check.lean}.

\begin{thebibliography}{9}

\bibitem{OsterwalderSeiler}
K.~Osterwalder and E.~Seiler,
\emph{Gauge field theories on a lattice},
Annals of Physics \textbf{110} (1978), 440--471.

\bibitem{GlimmJaffe}
J.~Glimm and A.~Jaffe,
\emph{Quantum Physics: A Functional Integral Point of View},
2nd ed., Springer, 1987.

\bibitem{ErikssonOSChain}
L.~Eriksson,
\emph{From the Gibbs Weight to the Spectral Gap: A Complete Machine-Checked
Osterwalder--Seiler Chain for the $\Z_2$ Lattice Gauge Chain}, viXra, 2026.

\bibitem{ErikssonQuotient}
L.~Eriksson,
\emph{The Quotient That Is Not the Identity: A Machine-Checked Degenerate
Reflection Pairing and Its Gelfand--Naimark--Segal Quotient}, viXra, 2026.

\bibitem{Lean4}
L.~de~Moura and S.~Ullrich,
\emph{The Lean 4 theorem prover and programming language},
CADE-28, Lecture Notes in Computer Science \textbf{12699} (2021), 625--635.

\bibitem{Mathlib}
The mathlib Community,
\emph{The Lean mathematical library},
CPP 2020, 367--381.

\bibitem{Clay}
A.~Jaffe and E.~Witten,
\emph{Quantum Yang--Mills theory},
Clay Mathematics Institute Millennium Prize Problem description, 2000.

\end{thebibliography}

\end{document}
